\documentclass{article}

\usepackage{amsmath,amssymb}
\usepackage{xurl}
\usepackage[hidelinks]{hyperref}
\setlength{\emergencystretch}{2em}

% Shared commands for notes in docs/paper-gaps/.

\DeclareUnicodeCharacter{2080}{\ifmmode{}_{0}\else\textsubscript{0}\fi}
\DeclareUnicodeCharacter{2081}{\ifmmode{}_{1}\else\textsubscript{1}\fi}
\DeclareUnicodeCharacter{2082}{\ifmmode{}_{2}\else\textsubscript{2}\fi}
\DeclareUnicodeCharacter{2099}{\ifmmode{}_{n}\else\textsubscript{n}\fi}

\newcommand{\Lean}{\textsc{Lean}}
\ExplSyntaxOn
\NewDocumentCommand{\leanid}{m}{
  \ifmmode
    \textsf{\detokenize{#1}}
  \else
    \group_begin:
    \urlstyle{sf}
    \tl_set:Nn \l_tmpa_tl {#1}
    \tl_replace_all:Nnn \l_tmpa_tl {\_} {_}
    \exp_args:NV \path \l_tmpa_tl
    \group_end:
  \fi
}
\ExplSyntaxOff
\providecommand{\tr}{\operatorname{tr}}
\newcommand{\tnleanIssueBase}{https://github.com/LionSR/TNLean/issues/}
\newcommand{\tnleanPrBase}{https://github.com/LionSR/TNLean/pull/}
\newcommand{\tnleanDocBase}{https://sirui-lu.com/TNLean/docs/}
\newcommand{\ghissue}[1]{\href{\tnleanIssueBase#1}{GitHub issue \##1}}
\newcommand{\ghpr}[1]{\href{\tnleanPrBase#1}{PR \##1}}
\newcommand{\doclink}[1]{\href{\tnleanDocBase#1.html}{\path{#1}}}

% Dirac notation and the identity operator, as in blueprint/src/macros/common.tex.
% \providecommand keeps note-local definitions authoritative.
\usepackage{dsfont}
\providecommand{\ket}[1]{|#1\rangle}
\providecommand{\bra}[1]{\langle #1|}
\providecommand{\braket}[2]{\langle #1 | #2 \rangle}
\providecommand{\ketbra}[2]{|#1\rangle\!\langle #2|}
\providecommand{\Id}{\mathds{1}}
\providecommand{\one}{\mathds{1}}
\providecommand{\C}{\mathbb{C}}
\providecommand{\MN}[1]{M_{#1}(\C)}
\providecommand{\MD}{M_D(\C)}

% External referencing for paper sources.  The build helper writes sanitized
% auxiliary files under build/paper-aux/ and excludes duplicated source labels.
\usepackage{xr}

% Import a generated paper auxiliary file with a caller-supplied label prefix.
% The candidate paths support builds from docs/paper-gaps/, the repository root,
% and build/paper-aux/ itself.
\newcommand{\importpaperaux}[2]{%
  \IfFileExists{../../build/paper-aux/#2.aux}{%
    \externaldocument[#1]{../../build/paper-aux/#2}%
  }{%
    \IfFileExists{build/paper-aux/#2.aux}{%
      \externaldocument[#1]{build/paper-aux/#2}%
    }{%
      \IfFileExists{#2.aux}{%
        \externaldocument[#1]{#2}%
      }{}%
    }%
  }%
}

% General external reference macros
\newcommand{\paperref}[2]{\ref{#1-#2}}
\newcommand{\papereqref}[2]{(\ref{#1-#2})}

% CPSV16 source (1606.00608 / MPDO-22-12-17-2.tex) external document and shortcuts
\importpaperaux{cpsv16-}{1606.00608/MPDO-22-12-17-2-tnlean}
\newcommand{\cpsvref}[1]{\paperref{cpsv16}{#1}}
\newcommand{\cpsveqref}[1]{\papereqref{cpsv16}{#1}}

% Verdict marker: \gapnote{<kind>}{<status>}. Kind is the policy
% classification (clarification, local-correction, scope-restriction,
% unfaithful, false-source, open-gap); status is open, wip (elimination
% underway), resolved, or historical. Machine-read by the site index;
% prints a verdict line.
\newcommand{\gapnote}[2]{\par\noindent\textbf{Verdict:} #1 (#2).\par}


\title{Per-Pair Support in the CPSV16 Figure 11 Decomposition}
\author{}
\date{2026-07-21}

\begin{document}
\maketitle
\gapnote{clarification}{resolved}

\section{The summed closed-chain identity}

Appendix~C.4 of Ref.~\cite{Cirac2016MPDO_arXiv} compares two closed-chain
expressions for the vertically read two-site tensor $B$. For a positive chain
length $L$, let $O_L$ denote contraction of a closed chain of length $L$. The
comparison is
\begin{align}
    O_L(B)
        &= \sum_\gamma \frac{m_\gamma^L}{n_\gamma}
            \tr(\nu_\gamma^L)O_L(M_\gamma)
    \label{eq:cpsv16_figure11_per_pair_support_global_closed_chain_sector_sum}\\
        &= \sum_{\alpha,\beta}
            \tr(\mu_\alpha^L)\tr(\mu_\beta^L)
            O_L(M_\alpha)O_L(M_\beta).
    \label{eq:cpsv16_figure11_per_pair_support_global_closed_chain_pair_sum}
\end{align}
The source gives this comparison in Appendix~C.4, lines~2011--2020 of
Ref.~\cite{Cirac2016MPDO_arXiv}. It then decomposes the products
$M_\alpha M_\beta$ into the normal sectors $M_\gamma$ shown in Figure~11,
with positive diagonal multiplicities $\chi_{\alpha,\beta,\gamma}$; see
Appendix~C.4, lines~2020--2029 of Ref.~\cite{Cirac2016MPDO_arXiv}.

Equations~(\ref{eq:cpsv16_figure11_per_pair_support_global_closed_chain_sector_sum})
and~(\ref{eq:cpsv16_figure11_per_pair_support_global_closed_chain_pair_sum})
give an identity only after summation over all ordered pairs
$(\alpha,\beta)$. They do not imply an identity for any fixed pair. More
generally, after repeated copies in $\mu_\alpha$ and $\mu_\beta$ have been
separated, an equality of the form
\begin{align}
    O_L(B)
        &= \sum_{p\in\mathcal P}O_L(C_p)
    \label{eq:cpsv16_figure11_per_pair_support_summed_copy_pairs}
\end{align}
does not determine any individual summand $O_L(C_p)$. A contribution may
vanish, and distinct decompositions of the total sum may distribute
contributions differently among the pair labels. Linear independence of the
vectors generated by the normal tensors separates the final normal label
$\gamma$ only after all pair contributions have been collected; it does not
recover the pair label $p$.

Consequently, equality of the global closed-chain sums cannot show that a
fixed copy pair has nonzero support or supply a positive normal corner for
that pair. Such a conclusion requires tensor-level information that retains
the pair coordinate before the closing trace.

\section{Tensor-level copy-pair corners}

Let $\mathcal A$ be the finite index set of retained copies, where each
$a\in\mathcal A$ records a normal label together with one copy of that label.
For weights $w_a$ and normal tensors $M_a$, define the retained one-site
vertical tensor by
\begin{align}
    A^i
        &= \bigoplus_{a\in\mathcal A}w_aM_a^i.
    \label{eq:cpsv16_figure11_per_pair_support_retained_one_site_tensor}
\end{align}
Let $\mathcal P=\mathcal A\times\mathcal A$ be the set of ordered copy pairs.
The product tensor $C$ has the canonical dependent direct-sum decomposition
\begin{align}
    C^u
        &= \bigoplus_{p\in\mathcal P}C_p^u,
    \label{eq:cpsv16_figure11_per_pair_support_pair_direct_sum}\\
    C_{(a,b)}^u
        &= w_aw_b(M_aM_b)^u.
    \label{eq:cpsv16_figure11_per_pair_support_pair_block_definition}
\end{align}
For $p\in\mathcal P$, let $E_p$ be the canonical inclusion of the bond space
of $C_p$ into the bond space of $C$. These inclusions satisfy
\begin{align}
    E_p^\dagger E_p
        &= I,
    \label{eq:cpsv16_figure11_per_pair_support_pair_inclusion_isometry}\\
    C^uE_p
        &= E_pC_p^u,
    \label{eq:cpsv16_figure11_per_pair_support_pair_right_intertwining}\\
    E_p^\dagger C^u
        &= C_p^uE_p^\dagger,
    \label{eq:cpsv16_figure11_per_pair_support_pair_left_intertwining}\\
    C^u
        &= \sum_{p\in\mathcal P}E_pC_p^uE_p^\dagger,
    \label{eq:cpsv16_figure11_per_pair_support_pair_corner_reconstruction}\\
    E_p^\dagger E_q
        &= 0
        \qquad (p\neq q).
    \label{eq:cpsv16_figure11_per_pair_support_distinct_pair_orthogonality}
\end{align}
These are tensor-letter identities, not consequences of the closed-chain
comparison in
Eq.~(\ref{eq:cpsv16_figure11_per_pair_support_global_closed_chain_pair_sum}).
The range projection $E_pE_p^\dagger$ therefore selects a fixed pair before
any word is evaluated or traced.

The vertically read one-site tensor $\widetilde M$ is related to the retained
tensor in
Eq.~(\ref{eq:cpsv16_figure11_per_pair_support_retained_one_site_tensor}) by an
exact coisometric reconstruction. There is a map $U$ such that
\begin{align}
    \widetilde M^i
        &= U^\dagger A^iU,
    \label{eq:cpsv16_figure11_per_pair_support_one_site_reconstruction}\\
    UU^\dagger
        &= I.
    \label{eq:cpsv16_figure11_per_pair_support_one_site_coisometry}
\end{align}
Squaring this identity and applying the standard product-coordinate
equivalences gives a coisometry $\widehat U$, obtained from $U\otimes U$, for
which
\begin{align}
    \widetilde{M^{[2]}}{}^u
        &= \widehat U^\dagger C^u\widehat U,
    \label{eq:cpsv16_figure11_per_pair_support_squared_reconstruction}\\
    \widehat U\widehat U^\dagger
        &= I.
    \label{eq:cpsv16_figure11_per_pair_support_squared_coisometry}
\end{align}
The squared reconstruction identifies the complete retained product tensor,
while the inclusions $E_p$ isolate its copy-pair corners. Together, these
identities provide the tensor-level support information absent from the
global closed-chain sum.

Every nonzero pair block $C_p$ has an exact spectral decomposition into
positive multiples of normal tensors:
\begin{align}
    C_p^u
        &= \sum_{k\in K_p}
            V_{p,k}(\lambda_{p,k}N_{p,k}^u)V_{p,k}^\dagger,
    \label{eq:cpsv16_figure11_per_pair_support_local_spectral_decomposition}\\
    \lambda_{p,k}
        &> 0
        \qquad (k\in K_p),
    \label{eq:cpsv16_figure11_per_pair_support_positive_spectral_weights}
\end{align}
where the $V_{p,k}$ are mutually orthogonal isometries. The composites
$E_pV_{p,k}$ retain both the pair label and the active normal corner.
Composing again with $\widehat U^\dagger$ places the same corner in the bond
space of the blocked vertical tensor. Flattening the dependent family
$\{(p,k):p\in\mathcal P,\ k\in K_p\}$ into a finite label set preserves all
support information.

The tensor-level reconstruction, pair-corner identities, and flattened
positive-length equality are formalized in
\doclink{TNLean/MPS/MPDO/VerticalProductFusionDecomposition.lean}.\footnote{The
principal declarations are
\leanid{MPOTensor.verticalTensor_blockTwo_squared_coisometry_reconstruction},
\leanid{MPOTensor.retainedVerticalProductTensor_eq_sum_pairBlocks},
\leanid{MPOTensor.RetainedProductSpectralFamily.retained_reconstruction_active},
and
\leanid{MPOTensor.RetainedProductSpectralFamily.flat_sameMPV}%
\textsubscript{2}\leanid{Pos}.}

\section{Zero copy pairs}

The active set $K_p$ in
Eq.~(\ref{eq:cpsv16_figure11_per_pair_support_local_spectral_decomposition})
may be empty. Two cases occur. First, the pair bond space may have dimension
zero. The canonical inclusion from this empty space remains an isometry
because its isometry equation is an equality of matrices on the empty index
type. Second, the pair bond space may be nonzero while every letter of $C_p$
vanishes. The entire pair space is then a zero corner, and exact
reconstruction is the empty sum
\begin{align}
    C_p^u
        &= 0
    \label{eq:cpsv16_figure11_per_pair_support_zero_pair_block}\\
        &= \sum_{k\in\varnothing}
            V_{p,k}(\lambda_{p,k}N_{p,k}^u)V_{p,k}^\dagger.
    \label{eq:cpsv16_figure11_per_pair_support_empty_pair_reconstruction}
\end{align}
No positive coefficient or artificial normal tensor should be assigned to
such a pair.

This convention agrees with the coisometric orientation of the Figure~11
fusion map in Ref.~\cite{Cirac2016MPDO_arXiv}. A zero pair has an empty active
target, and the unique map onto that target is a coisometry. Requiring an
isometry on the entire pair bond space would impose a full-support hypothesis
and exclude the elementary zero cross-products described in
\path{docs/paper-gaps/cpsv16_figure11_fusion_coisometry.tex}.

For the source multiplicities, an absent normal label is represented by an
empty $\chi_{\alpha,\beta,\gamma}$ fiber, not by a spurious positive diagonal
entry. Its contribution to every positive-length power sum is zero. For a
fixed ordered pair $(\alpha,\beta)$, the Figure~11 decomposition in
Appendix~C.4, lines~2020--2029 of Ref.~\cite{Cirac2016MPDO_arXiv}, is
\begin{align}
    U_{\alpha,\beta}(M_\alpha M_\beta)^iU_{\alpha,\beta}^\dagger
        &= \bigoplus_\gamma
            \chi_{\alpha,\beta,\gamma}\otimes
            Z_{\alpha,\beta,\gamma}M_\gamma^i
            Z_{\alpha,\beta,\gamma}^{-1}.
    \label{eq:cpsv16_figure11_per_pair_support_figure11_pair_decomposition}
\end{align}
Here $\gamma$ ranges over the normal sectors that occur for the chosen pair.
If $M_\alpha M_\beta$ vanishes, the right-hand side is the empty direct sum.
Every nonzero normal corner of a fixed pair is gauge-phase equivalent to a
sector $M_\gamma$. Absent sectors are omitted rather than represented by zero
coefficients. Equality of the global closed-chain sums over all
$(\alpha,\beta)$ cannot strengthen this statement to nonempty support for
every fixed pair.

\section{Verdict}

The global closed-chain comparison in Appendix~C.4 of
Ref.~\cite{Cirac2016MPDO_arXiv} determines the sum of all copy-pair
contributions but does not isolate a fixed pair. Per-pair support instead
follows from exact squared reconstruction and the canonical orthogonal
corners of the retained product tensor. Spectral decomposition is performed
separately on each corner, with an empty active family when the corner is
zero. This construction assumes neither full support nor normal sectors with
identically zero coefficients.

\bibliographystyle{alpha}
\bibliography{references}

\end{document}
